% unit: noether.p08.title.001
\section*{8. 关于任意多个基本形式的系统的不变量的整有理表示}
\begin{center}
Math. Ann. 77 (1916), pp. 93--102
\end{center}

% unit: noether.p08.par.001
在他为 Schwarz 纪念文集所写的《任意多个基本形式的系统的不变量》一文的结尾，D. Hilbert 提出猜想：这些不变量可以由系统 $(J,PJ)$ 的有限多个不变量整有理地表示，其中 $J$ 表示那些线性无关的基本形式的完整不变量系统，而不变量 $PJ$ 则由 $J$ 通过极化过程产生。

% unit: noether.p08.par.002
下面我在 I 中说明，这个猜想确实成立，而且它只是如下归约定理的简单推论：任意多个、每个含 $n$ 个变量的行的任一形式，都可表示为若干极化形式的和；这些极化形式只含有 $n$ 个固定行，并且它们本身由原始形式通过极化过程导出。这个归约定理是 F. Mertens 首先表述的一个特殊情形，属于 A. Capelli、F. Mertens 和 J. Deruyts 给出的 Clebsch--Gordan 级数展开推广；该推广适用于任意多个、每个含 $n$ 个变量的行的形式，并通过对微分过程的形式变换而证明。\footnote{F. Mertens: ``Über eine Formel der Determinantentheorie.'' (公式 5), Sitzb. d. Ak. d. Wiss. Wien, vol. 91, II. Abt. (1885), 以及更详尽的（含应用）``Über ganze Funktionen von $m$ Systemen von je $n$ Unbestimmten''. Monatsh. f. Math. u. Phys. 第 4 年 (1893)。关于 Capelli 自 1882 年以来的若干证明，参见例如 Math. Ann. 37，或其影印讲义 ``Lezioni sulla teoria delle forme algebriche'' (1902)。J. Deruyts: \emph{Essai d'une théorie générale des formes algébriques}, Mém. d. l. Soc. Roy. des Sciences de Liège (2) 17 (1892).}

% unit: noether.p08.par.003
在 II 中，我还给出归约定理的一个更概念化的证明，把这个问题看作线性形式族的等价问题。这种观点，以及所用论证中的一个实质部分，取自 E. Fischer 的论文“论代数模系统”\footnote{E. Fischer: Über algebraische Modulsysteme und lineare homogene partielle Differentialgleichungen mit konstanten Koeffizienten, \S\ 1--4, J. f. M. 140 (1911).}以及与之相连的 E. Fischer 未发表定理；他友好地允许我使用这些结果。

% unit: noether.p08.par.004
最后，在 III 中我还说明，对应的论证如何导向最一般的级数展开。

% unit: noether.p08.sec.001
\subsection*{I.}
% unit: noether.p08.par.005
设 $\theta$ 是在每一个 $N>n$ 行
% unit: noether.p08.eq.001
\[
  A_1,A_2,\ldots,A_N,
\]
中都齐次的形式，其中一般地，行 $A_k$ 由 $n$ 个元素
% unit: noether.p08.eq.002
\[
  A_k^{(1)},A_k^{(2)},\ldots,A_k^{(n)}
\]
组成。$\theta$ 的系数可以是不定元，也可以是某个给定有理性域中的量。又令 $P$ 表示极化过程
% unit: noether.p08.eq.003
\[
  P_{hk}=\left(A_h\frac{\partial}{\partial A_k}\right)
  =A_h^{(1)}\frac{\partial}{\partial A_k^{(1)}}+
   A_h^{(2)}\frac{\partial}{\partial A_k^{(2)}}+
   \cdots+
   A_h^{(n)}\frac{\partial}{\partial A_k^{(n)}}
\]
的一个整有理函数，其系数为有理整数；其中乘积 $P_{hk}P_{ij}\theta$ 应理解为把 $P_{hk}$ 施加于 $P_{ij}\theta$。

% unit: noether.p08.par.006
刚才提到的归约定理由恒等式
% unit: noether.p08.eq.004
\begin{equation}
  \theta=\sum PZ,
\tag{1}
\end{equation}
给出，其中形式 $Z$ 只含有行
% unit: noether.p08.eq.005
\[
  A_1,A_2,\ldots,A_n
\]
并且由 $\theta$ 通过极化过程 $P$ 导出。

% unit: noether.p08.par.007
现在考虑基本形式系统 $(F')$，它由 $N$ 个同阶且变量数相同的形式
% unit: noether.p08.eq.006
\[
  F_1,F_2,\ldots,F_N
\]
组成，其系数分别取作 $N$ 个行
% unit: noether.p08.eq.007
\[
  A_1,A_2,\ldots,A_N.
\]
令 $(J)$ 表示 $F_1,\ldots,F_n$ 的完整系统，它由有限多个不变量组成，并且使每个 $F_1,\ldots,F_n$ 的不变量都可由 $(J)$ 中的不变量整有理地表示。要证明的 Hilbert 猜想便是：对于 $F_1,\ldots,F_N$ 的同时不变量 $S$，一个完整系统由有限多个不变量 $(J,PJ)$ 给出，其中 $PJ$ 表示所有由 $J$ 通过极化过程 $P$ 导出的不变量。

% unit: noether.p08.par.008
事实上，$(F')$ 的每一个带有有理整数系数的同时不变量 $S$，都是在 $A_1,\ldots,A_N$ 的每一个行中齐次的形式，且系数为有理整数，因此可对它应用恒等式 (1)。由于对 $S$ 施加极化过程 $P$ 已知仍产生不变量，所以形式 $Z$ 也都是不变量；并且由于它们只含有行 $A_1,\ldots,A_n$，它们是 $F_1,\ldots,F_n$ 的同时不变量；于是按假设，它们是 $(J)$ 中不变量的整有理函数。因此恒等式 (1) 变成
% unit: noether.p08.eq.008
\[
  S=\sum PY,
\]
其中 $Y$ 表示 $(J)$ 中不变量的幂积。可是由 $P$ 的定义可知，实际在 $Y$ 上施行过程 $P$，就是有限次相继施行简单的极化操作 $P_{hk}$；按乘积的微分规则，这导致 $(J)$ 与 $(PJ)$ 中不变量的乘积之和。对来自 $(J)$ 与 $(PJ)$ 的乘积施加 $P_{hk}$ 亦然；因此 $PY$ 也只给出 $(J,PJ)$ 的整有理组合。\emph{由此证明了所有不变量 $S$ 都可由 $(J,PJ)$ 整有理地表示。}

% unit: noether.p08.sec.002
